\chapter{مفاهیم پایه و مرور کارهای پیشین}
\label{ch:literature}

این فصل زمینه‌ی لازم برای فهم روش پیشنهادی را فراهم می‌کند. ساختار آن از عام به خاص است: نخست مفاهیم پایه‌ی یادگیری عمیق و اجزایی که در آزمایش‌ها به‌کار می‌روند مرور می‌شوند؛ سپس یادگیری پیوسته و خانواده‌های راه‌حل آن؛ آنگاه یادگیری‌زدایی ماشینی و تلاقی آن با یادگیری پیوسته؛ پس از آن یادگیری کم‌پیراسنجه؛ و در پایان روش‌شناسی ارزیابی فراموشی و حریم خصوصی. فصل با یک جدول مقایسه‌ی مفهومی و بیان صریح خلأ پژوهشی بسته می‌شود. برای مرورهای جامع یادگیری پیوسته می‌توان به \cite{parisi2019continual,delange2022survey,vandeven2022three,wang2023continualsurvey} و برای مرورهای یادگیری‌زدایی به \cite{nguyen2022unlearningsurvey,xu2023unlearningsurvey} رجوع کرد.

\section{مفاهیم پایه}
\label{sec:lit-basics}

\subsection{یادگیری بانظارت و بهینه‌سازی مبتنی بر شیب}
\label{subsec:lit-sgd}

یک شبکه‌ی طبقه‌بند را با $f(\cdot\,;\theta)$ نشان می‌دهیم که ورودی $x$ را به برداری از {امتیازهای خام\en{Logits}} می‌نگارد. آموزش بانظارت، کمینه‌سازی {مخاطره‌ی تجربی\en{Empirical Risk}} است:
\begin{equation}
	\label{eq:erm}
	\widehat{R}(\theta)=\frac{1}{|D|}\sum_{(x,y)\in D}\ell\big(f(x;\theta),y\big),
\end{equation}
که در آن $\ell$ معمولا زیان {بی‌نظمی متقاطع\en{Cross-Entropy}} است. کمینه‌سازی با {شیب کاهشی تصادفی\en{Stochastic Gradient Descent (SGD)}} \cite{robbins1951sgd} و به‌روزرسانی $\theta\!\leftarrow\!\theta-\eta\nabla_\theta\widehat{R}_{B}(\theta)$ روی دسته‌های کوچک $B$ انجام می‌شود. همه‌ی آزمایش‌های این پایان‌نامه از \lr{SGD} با تکانه استفاده می‌کنند؛ دلیل این انتخاب، حذف اثر تطبیقی بهینه‌سازهای مبتنی بر ممان دوم است، چون در حالت فراموشی، حالت انباشته‌ی بهینه‌ساز خودش می‌تواند اثری از وظیفه‌ی حذف‌شده حمل کند.

بردار شیب در این پایان‌نامه دو نقش دارد و تفکیک این دو نقش مهم است. نقش اول، نقش متعارف آن به‌عنوان جهت به‌روزرسانی است. نقش دوم — که سازوکار پیشنهادی بر آن استوار است — استفاده از \emph{بزرگی} مؤلفه‌های شیب به‌عنوان یک سنجه‌ی محلی از {اهمیت پیراسنجه\en{Parameter Importance}} است: اگر $|\partial_i \mathcal{L}_{\mathrm{retain}}|$ بزرگ باشد، جابه‌جایی کوچک در درایه‌ی $i$ زیان وظایف نگه‌داری‌شده را به‌طور محسوس تغییر می‌دهد.

\subsection{سنجه‌های اهمیت پیراسنجه}
\label{subsec:lit-importance}

پژوهش‌های یادگیری پیوسته چند سنجه‌ی اهمیت پیشنهاد کرده‌اند و انتخاب میان آن‌ها به موازنه‌ی دقت و هزینه بستگی دارد. {ماتریس اطلاعات فیشر\en{Fisher Information Matrix}} در \lr{EWC} \cite{kirkpatrick2017overcoming} از میانگین مربع شیب لگاریتم درست‌نمایی برآورد می‌شود و مبنای نظری روشنی دارد، اما محاسبه‌ی آن به یک گذر کامل روی داده‌ها نیاز دارد. {هوشمندی سیناپسی\en{Synaptic Intelligence}} \cite{zenke2017continual} اهمیت را در طول آموزش انباشته می‌کند و {سیناپس‌های آگاه از حافظه\en{Memory Aware Synapses}} \cite{aljundi2018memory} از حساسیت خروجی به تغییر وزن‌ها استفاده می‌کنند. {توابع نفوذ\en{Influence Functions}} \cite{koh2017influence} اثر یک نمونه‌ی آموزشی بر پیراسنجه‌ها را تقریب می‌زنند و پیوند مفهومی نزدیکی با یادگیری‌زدایی دارند.

در این پایان‌نامه، سنجه‌ی پیش‌فرض بزرگی شیب روی حافظه‌ی بازپخش وظایف نگه‌داری‌شده است. دلیل این انتخاب صرفا عمل‌گرایانه است: یک گذر پس‌رو روی حافظه‌ی محدود کافی است و هزینه‌ی آن نسبت به حلقه‌ی ترمیم ناچیز است. در بخش \ref{subsec:res-conflict} یک جانشین مبتنی بر {انرژی تعارض شیب‌ها\en{Gradient-Conflict Energy}} نیز به‌طور تجربی سنجیده می‌شود که نتیجه‌ی آن وابسته به مجموعه‌داده است.

\subsection{معماری‌های استفاده‌شده}
\label{subsec:lit-arch}

آزمایش‌ها بر دو خانواده‌ی معماری استوارند. خانواده‌ی نخست {شبکه‌های باقی‌مانده\en{Residual Networks}} یا \lr{ResNet} \cite{he2016deep} است. اتصال باقی‌مانده‌ی $y=\mathcal{F}(x)+x$ مسیری کوتاه برای جریان شیب ایجاد می‌کند و آموزش شبکه‌های عمیق را پایدار می‌سازد. در این کار \lr{ResNet-18} برای \lr{CIFAR-10} و \lr{TinyImageNet} و اجراهای پیش‌آموزش‌دیده، و \lr{ResNet-34} برای اجراهای اصلی \lr{CIFAR-100} به‌کار می‌رود. خانواده‌ی دوم {مبدل بینایی\en{Vision Transformer}} \cite{dosovitskiy2021vit} است که تصویر را به وصله‌ها تقسیم و با لایه‌های توجه پردازش می‌کند؛ از این خانواده یک نسخه‌ی کوچک به‌عنوان آزمون انتقال میان‌معماری استفاده می‌شود.

مقدارگذاری اولیه نیز در جای‌جای روش نقش دارد. بازنشانی درایه‌های وظیفه‌ی هدف از توزیع {کایمینگ\en{Kaiming Initialization}} \cite{he2015delving} نمونه می‌گیرد؛ همین انتخاب در لایه‌های تطبیق هم برای نگاشت کاهش ابعاد تکرار می‌شود، در حالی که نگاشت افزایش ابعاد با صفر مقدار می‌گیرد تا لایه‌ی تطبیق در لحظه‌ی تولد، تابع همانی باشد.

\subsection{حافظه‌ی محدود و بازپخش}
\label{subsec:lit-memory}

فرض دسترسی‌نداشتن به داده‌ی وظایف گذشته در عمل بسیار سخت‌گیرانه است. راه‌حل رایج، نگه‌داشتن یک {حافظه‌ی رویدادی\en{Episodic Memory}} کوچک است. در این کار بودجه‌ی بازپخش $500$ نمونه است و — نکته‌ی مهم — این بودجه از پیش میان وظایف افراز می‌شود و پس از یک درخواست حذف، سهم آزادشده به وظایف دیگر بازتوزیع نمی‌شود. دلیل این محدودیت، جلوگیری از یک اثر مداخله‌گر است: اگر بودجه بازتوزیع شود، بخشی از بهبود مشاهده‌شده پس از فراموشی می‌تواند صرفا ناشی از بزرگ‌ترشدن حافظه‌ی وظایف باقی‌مانده باشد و نه ناشی از سازوکار فراموشی.

\section{یادگیری پیوسته}
\label{sec:lit-cl}

\subsection{سه حالت یادگیری}
\label{subsec:lit-scenarios}

پیکربندی‌های یادگیری پیوسته را می‌توان بر اساس اطلاعات موجود در استنتاج به سه حالت یادگیری تقسیم کرد \cite{vandeven2022three}: {افزایشی‌وظیفه\en{Task-Incremental}} که شناسه‌ی وظیفه در استنتاج در دسترس است؛ {افزایشی‌دامنه\en{Domain-Incremental}} که مجموعه‌ی برچسب‌ها ثابت است و توزیع ورودی تغییر می‌کند؛ و {افزایشی‌کلاس\en{Class-Incremental}} که مدل باید بدون شناسه، میان همه‌ی کلاس‌های دیده‌شده تمایز بگذارد. این پایان‌نامه در حالت نخست کار می‌کند. این انتخاب یک محدودیت واقعی است، اما برای مسئله‌ی حذف در سطح وظیفه ضروری است: مفهوم «فراموش‌کردن وظیفه‌ی $f$» تنها زمانی خوش‌تعریف است که مرزهای وظیفه در سطح رابط سامانه شناخته باشند.

تفاوت این سه حالت یادگیری را می‌توان دقیق‌تر بیان کرد. فرض کنید مدل تابعی از ورودی $x$ و شناسه‌ی وظیفه $t$ باشد. در حالت افزایشی‌وظیفه، مدل $p(y \mid x, t)$ را می‌آموزد و در استنتاج هر دو ورودی در دسترس‌اند؛ عملا این یعنی هر وظیفه می‌تواند سر طبقه‌بند خودش را داشته باشد و مسئله‌ی تمایز میان کلاس‌های وظایف مختلف اصلا مطرح نمی‌شود. در حالت افزایشی‌دامنه، فضای برچسب ثابت است و مدل همچنان $p(y \mid x)$ را می‌آموزد، اما توزیع ورودی در طول زمان جابه‌جا می‌شود. دشوارترین حالت، حالت افزایشی‌کلاس است: مدل باید $p(y \mid x)$ را روی \emph{اجتماع} همه‌ی کلاس‌های دیده‌شده بیاموزد بدون آنکه بداند نمونه از کدام وظیفه آمده است، و بنابراین باید مرزهایی میان کلاس‌هایی بکشد که هرگز در یک دسته‌ی آموزشی کنار هم ظاهر نشده‌اند.

این تفکیک برای مسئله‌ی حذف پیامد مستقیم دارد. واحد حذف در این پایان‌نامه یک وظیفه‌ی کامل است، و «وظیفه» تنها در حالت نخست یک شیء صریح در رابط سامانه است. در حالت افزایشی‌کلاس، درخواست حذف باید در سطح کلاس یا رکورد بیان شود، که مسئله‌ی متفاوتی است: در آنجا نمی‌توان یک سر طبقه‌بند یا یک زیرشبکه‌ی مشخص را کنار گذاشت، چون دانش کلاس حذف‌شونده در همان مرزهایی پخش است که کلاس‌های باقی‌مانده را از هم جدا می‌کنند. بنابراین انتخاب حالت افزایشی‌وظیفه در این کار یک ساده‌سازی دلبخواهی نیست، بلکه شرط خوش‌تعریفی صورت‌مسئله است؛ در عوض، تعمیم روش به دو حالت دیگر در بخش \ref{sec:concl-limitations} به‌عنوان محدودیت ثبت می‌شود.

\subsection{روش‌های مبتنی بر تنظیم}
\label{subsec:lit-reg}

نخستین خانواده‌ی راه‌حل، جابه‌جایی پیراسنجه‌های مهم را جریمه می‌کند. \lr{EWC} \cite{kirkpatrick2017overcoming} جمله‌ی $\sum_i F_i(\theta_i-\theta_i^{*})^2$ را با وزن‌های فیشر اضافه می‌کند؛ \lr{SI} \cite{zenke2017continual} و \lr{MAS} \cite{aljundi2018memory} همان قالب را با سنجه‌های اهمیت متفاوت پیاده می‌کنند. {یادگیری بدون فراموشی\en{Learning without Forgetting}} یا \lr{LwF} \cite{li2017learning} به‌جای جریمه‌ی وزن، خروجی مدل جدید را به خروجی مدل قدیم نزدیک نگه می‌دارد و در واقع یک {تقطیر دانش\en{Knowledge Distillation}} \cite{hinton2015distilling} خودارجاع است.

شهود مشترک این خانواده ساده است: اگر بدانیم کدام پیراسنجه‌ها برای وظایف گذشته حیاتی‌اند، می‌توانیم به‌جای نگه‌داشتن داده، خود پیراسنجه‌ها را نگه داریم. اختلاف روش‌ها در پاسخ به این پرسش است که «اهمیت» چگونه اندازه‌گیری شود. \lr{EWC} از قطر ماتریس اطلاعات فیشر استفاده می‌کند، که برابر امید مربع شیب لگاریتم درست‌نمایی است و بنابراین حساسیت \emph{محلی} خروجی مدل به هر پیراسنجه را می‌سنجد. \lr{SI} همان کمیت را به‌صورت برخط و در طول مسیر بهینه‌سازی انباشته می‌کند و به یک گذر جداگانه نیاز ندارد. \lr{LwF} مسیر متفاوتی می‌رود و به‌جای پیراسنجه‌ها، \emph{رفتار} مدل را مقید می‌کند: خروجی نرم مدل پیشین روی داده‌ی جدید به‌عنوان هدف تقطیر به‌کار می‌رود، پس هیچ داده‌ای از وظایف گذشته لازم نیست.

دو محدودیت مشترک این خانواده در ادبیات ثبت شده است. نخست، جریمه سراسری است و روی همه‌ی پیراسنجه‌ها اعمال می‌شود؛ با افزایش تعداد وظایف، جمله‌های جریمه روی هم انباشته می‌شوند و ظرفیت یادگیری وظیفه‌ی تازه را می‌بلعند. دوم، تقریب قطری فیشر همبستگی میان پیراسنجه‌ها را نادیده می‌گیرد، در حالی که آسیب واقعی اغلب از حرکت \emph{هم‌زمان} گروهی از وزن‌ها می‌آید.

پیوند این خانواده با کار حاضر نزدیک است، اما جهت آن معکوس. جمله‌ی لنگر $\ell_2$ در روش پیشنهادی از نظر شکل ریاضی همان قالب \lr{EWC} است؛ تفاوت در این است که در \lr{EWC} هدف، محافظت در برابر \emph{یادگیری} وظیفه‌ی بعدی است و مجموعه‌ی محافظت‌شده کل پیراسنجه‌های مهم است، در حالی که در این کار هدف، محافظت در برابر \emph{فراموش‌کردن} است و مجموعه‌ی محافظت‌شده به هم‌پوشانی ساختاری میان زیرشبکه‌ی هدف و زیرشبکه‌های فعال محدود می‌شود. همچنین یک نکته‌ی تجربی مهم در فصل \ref{ch:results} روشن می‌شود: \lr{EWC} و \lr{LwF} در معیار بیشترین افت عالی‌اند اما در سرکوب وظیفه‌ی هدف شکست می‌خورند، چون سازوکار آن‌ها اساسا برای حفظ ساخته شده است، نه برای حذف.

\subsection{روش‌های مبتنی بر بازپخش}
\label{subsec:lit-replay}

خانواده‌ی دوم، نمونه‌های گذشته را در حافظه نگه می‌دارد و در آموزش وظیفه‌ی جدید بازپخش می‌کند. {بازپخش تجربه\en{Experience Replay}} ساده‌ترین شکل آن است و \lr{iCaRL} \cite{rebuffi2017icarl} آن را با نمونه‌های شاخص و طبقه‌بندی نزدیک‌ترین میانگین ترکیب می‌کند. {بازپخش تجربه‌ی تاریک\en{Dark Experience Replay}} یا \lr{DER++} \cite{buzzega2020dark} علاوه بر برچسب‌ها، امتیازهای خام ذخیره‌شده را نیز تطبیق می‌دهد و به همین دلیل اطلاعات غنی‌تری از رفتار گذشته‌ی مدل حمل می‌کند. جایگزین مولد این خانواده، بازپخش نمونه‌های ساختگی است \cite{shin2017deepgenerative}.

زیان ترمیمی روش پیشنهادی مستقیما از این خط می‌آید: جمله‌ی $\mathcal{L}_{DER}$ همان تطبیق امتیازهای خام ذخیره‌شده است. دلیل این انتخاب آن است که پس از یک درخواست حذف، هدف ما بازگرداندن \emph{رفتار پیشین} مدل روی وظایف فعال است و امتیازهای خام ذخیره‌شده دقیقا همان رفتار را ثبت کرده‌اند.



اندازه‌ی حافظه در این خانواده یکی از عوامل تعیین‌کننده در طراحی است. حافظه‌ی بازپخش نسخه‌ای از داده‌های کاربر را نگه می‌دارد؛ بنابراین اگر پس از درخواست حذف یک وظیفه، نمونه‌های بازپخش آن همچنان باقی بمانند، داده‌های آن وظیفه هنوز در سامانه حضور دارند و ادعای حذف پذیرفتنی نیست. به همین دلیل، حذف نمونه‌های بازپخش وظیفه‌ی هدف بخشی از عملیات حذف در این پایان‌نامه است و فقط یک جزئیات پیاده‌سازی نیست. بودجه‌ی حافظه پس از حذف دوباره میان وظیفه‌ها توزیع نمی‌شود تا همه‌ی روش‌ها با بودجه‌ی یکسان مقایسه شوند.

\subsection{تصویر شیب}
\label{subsec:lit-projection}

{حافظه‌ی مبتنی بر شیب رویدادی\en{Gradient Episodic Memory}} یا \lr{GEM} \cite{lopezpaz2017gem} و نسخه‌ی میانگین‌گیری‌شده‌ی آن \lr{A-GEM} \cite{chaudhry2019agem} به‌روزرسانی وظیفه‌ی جدید را طوری تصویر می‌کنند که ضرب داخلی آن با شیب وظایف گذشته منفی نشود. این خانواده از نظر ایده نزدیک‌ترین کار به سازوکار این پایان‌نامه است: در هر دو، شیب وظایف نگه‌داری‌شده به‌روزرسانی را مقید می‌کند. تفاوت کلیدی در محل اعمال قید است. \lr{GEM} یک تصویر \emph{سراسری} در فضای شیب انجام می‌دهد و همه‌ی درایه‌ها را به‌طور یکنواخت تحت تأثیر قرار می‌دهد؛ روش این پایان‌نامه یک ناحیه‌ی \emph{انتخاب‌شده} از درایه‌های مشترک را شناسایی و در همان‌جا مهار می‌کند. این تفاوت هم مفهومی است و هم عملی: انتخاب ناحیه‌ای اجازه می‌دهد بقیه‌ی به‌روزرسانی آزادانه در خدمت فراموشی باقی بماند.

\subsection{جداسازی پیراسنجه و زیرشبکه‌های تنک}
\label{subsec:lit-isolation}

خانواده‌ی چهارم، پیراسنجه‌ها را به‌ازای هر وظیفه جدا می‌کند. \lr{PackNet} \cite{mallya2018packnet} پس از آموزش هر وظیفه هرس می‌کند و وزن‌های باقی‌مانده را قفل می‌کند؛ {نقاب توجه سخت\en{Hard Attention to the Task}} یا \lr{HAT} \cite{serra2018hat} نقاب‌های واحدمحور یاد می‌گیرد؛ و \lr{SupSup} \cite{wortsman2020supsup} تنها نقاب‌های دودویی را روی یک بدنه‌ی تصادفی و ثابت می‌آموزد. {شبکه‌های پیش‌رونده\en{Progressive Networks}} \cite{rusu2016progressive} در انتهای دیگر این طیف قرار دارند و ظرفیت را به‌ازای هر وظیفه رشد می‌دهند.

مبنای نظری این خانواده به {فرضیه‌ی بلیت بخت‌آزمایی\en{Lottery Ticket Hypothesis}} \cite{frankle2019lottery} و به نتیجه‌ی مکمل آن \cite{ramanujan2020hidden} می‌رسد: زیرشبکه‌های قابل‌آموزش با کارایی خوب، درون شبکه‌های چگال و حتی درون شبکه‌های تصادفی وجود دارند. {یادگیری پیوسته‌ی زیرشبکه‌ی برنده\en{Winning Subnetwork Continual Learning}} یا \lr{WSN} \cite{kang2022wsn} این ایده را به یادگیری پیوسته می‌آورد و برای هر وظیفه یک نقاب دودویی می‌آموزد. \lr{PALL} نقاب‌های هر وظیفه‌ی خود را از همین خط به ارث می‌برد و همین ارث‌بری است که مفهوم \emph{هم‌پوشانی ساختاری} میان دو وظیفه را خوش‌تعریف می‌کند: هم‌پوشانی، اشتراک مجموعه‌ای دو نقاب است و لازم نیست با یک سنجه‌ی پیوسته تقریب زده شود. این نکته پایه‌ی مستقیم بخش \ref{sec:prop-epall} است.



چارچوب \lr{PALL} \cite{ozdenizci2025pall} همین ایده را برای حذف به‌کار می‌گیرد. برای هر وظیفه، یک نقاب دودویی تنک روی وزن‌های شبکه آموخته و ذخیره می‌شود تا هر وظیفه فقط از زیرشبکه‌ی خود استفاده کند. در زمان حذف، نقاب وظیفه‌ی هدف بازنشانی، سر طبقه‌بند مربوط به آن پاک و نمونه‌های بازپخش آن وظیفه از حافظه حذف می‌شوند. مزیت این طراحی آن است که هزینه‌ی حذف به اندازه‌ی داده وابسته نیست و برخلاف روش‌های مبتنی بر صعود از روی شیب، به بهینه‌سازی روی داده‌های در حال حذف نیاز ندارد.

نقطه‌ی ضعف روش و انگیزه‌ی این پایان‌نامه در همین بخش آشکار می‌شود. نقاب‌های وظیفه‌های مختلف کاملا جدا نیستند. در یک شبکه با تنکی ثابت، هرچه تعداد وظیفه‌ها بیشتر شود، اشتراک نقاب‌ها نیز افزایش می‌یابد. در نتیجه، بازنشانی زیرشبکه‌ی هدف ممکن است وزن‌هایی را تغییر دهد که وظیفه‌های فعال هنوز به آن‌ها وابسته‌اند. \lr{PALL} برای جبران این آسیب، پس از بازنشانی یک مرحله‌ی ترمیم روی داده‌های بازپخش وظیفه‌های باقی‌مانده اجرا می‌کند؛ اما در این مرحله میان وزن‌های مشترک و غیرمشترک تفاوتی گذاشته نمی‌شود و سازوکاری برای پیشگیری از آسیب وجود ندارد. پژوهش حاضر همین کاستی را هدف قرار می‌دهد.

\section{یادگیری‌زدایی ماشینی}
\label{sec:lit-unlearning}

\subsection{فراموشی دقیق و بازآموزی}
\label{subsec:lit-exact}

مرجع مفهومی یادگیری‌زدایی، بازآموزی کامل مدل روی داده‌ی باقی‌مانده است. چون این کار پرهزینه است، \lr{SISA} \cite{bourtoule2021machine} داده را به بخش‌های جدا افراز می‌کند، یک عضو گروه را روی هر بخش آموزش می‌دهد و در زمان حذف تنها عضو آسیب‌دیده را بازآموزی می‌کند. {آموزش فراموش‌کار\en{Amnesiac Training}} \cite{graves2021amnesiac} به‌روزرسانی‌های هر دسته را ذخیره می‌کند تا بعدا کم شود. کارهای دیگری برای مدل‌های ساده مانند خوشه‌بندی \lr{k-means} حذف دقیق ارائه کرده‌اند \cite{ginart2019making}.

معنای دقیق «دقیق» در این خانواده اهمیت دارد. یک سازوکار حذف را دقیق می‌نامیم هرگاه توزیع مدل خروجی پس از حذف نمونه‌ی $z$، با توزیع مدلی که از ابتدا بدون $z$ آموزش دیده است \emph{یکسان} باشد. این تعریف بر توزیع مدل‌ها بنا شده و نه بر رفتار یک مدل منفرد؛ به همین دلیل صرف افت دقت روی داده‌ی حذف‌شده، شاهدی برای حذف دقیق نیست. بازآموزی کامل به‌طور بدیهی این تعریف را برآورده می‌کند و به همین دلیل مرجع مفهومی این حوزه است.

هزینه‌ی این مرجع یکی از مسائل اصلی این حوزه است. بازآموزی کامل به نگه‌داری دائمی داده‌های خام و تکرار تمام هزینه‌ی آموزش در هر درخواست نیاز دارد. در محیط پیوسته که درخواست‌ها به‌صورت جریانی می‌رسند، این هزینه برای هر درخواست دوباره پرداخت می‌شود. رویکردهای افرازی مانند \lr{SISA} هزینه را کاهش می‌دهند، اما دو پیامد دارند: کاهش دقت، چون هر عضو گروه فقط بخشی از داده‌ها را می‌بیند، و افزایش حافظه‌ی موردنیاز، چون چند مدل باید هم‌زمان نگه‌داری شوند. همین موازنه در مقایسه با روش جداسازی کامل نیز بررسی می‌شود.

\subsection{فراموشی تقریبی}
\label{subsec:lit-approx}

خانواده‌ی تقریبی به‌جای بازآموزی، وزن‌ها را مستقیما ویرایش می‌کند. \cite{golatkar2020forgetting} یک به‌روزرسانی مبتنی بر فیشر برای پاک‌کردن اطلاعات پیشنهاد می‌دهد؛ \cite{izzo2021approximate} روش‌های حذف تقریبی برای مدل‌های خطی ارائه می‌کند؛ {فراموشی مرزی\en{Boundary Unlearning}} \cite{chen2023boundary} مرزهای تصمیم خروجی را جابه‌جا می‌کند تا سرعت به‌دست آورد؛ و \cite{chundawat2023bad} از یک جفت معلم خوب و بد برای هدایت فراموشی استفاده می‌کند. دو کار اخیرتر به این پایان‌نامه نزدیک‌ترند و در آزمایش‌ها به‌عنوان خط پایه به‌کار می‌روند: {یادگیری‌زدایی برجستگی\en{Saliency Unlearning}} یا \lr{SalUn} \cite{fan2024salun} به‌روزرسانی را به وزن‌های پرنفوذ محدود می‌کند، و {میراکردن انتخابی سیناپسی\en{Selective Synaptic Dampening}} یا \lr{SSD} \cite{foster2024ssd} پیراسنجه‌ها را بر پایه‌ی اهمیت فیشری میرا می‌کند و برای این کار به هیچ گام آموزشی نیاز ندارد.

\lr{SalUn} از منظر این پایان‌نامه اهمیت مفهومی خاصی دارد، چون آن هم به‌روزرسانی را روی یک زیرمجموعه‌ی انتخاب‌شده متمرکز می‌کند. اما جهت انتخاب متفاوت است: \lr{SalUn} درایه‌های پرنفوذ را برای \emph{تغییردادن} انتخاب می‌کند، در حالی که روش این پایان‌نامه درایه‌های حساس به نگه‌داری را برای \emph{تغییرندادن} انتخاب می‌کند.

\subsection{تضمین‌های صوری و مرزهای آن}
\label{subsec:lit-certified}

یک شاخه‌ی نظری برای مدل‌های محدب، تضمین حذف ارائه می‌دهد. پژوهش \cite{guo2020certified} مفهوم {حذف داده‌ی گواهی‌شده\en{Certified Data Removal}} را با تعریفی شبیه به {حریم خصوصی تفاضلی\en{Differential Privacy}} \cite{dwork2014algorithmic} معرفی می‌کند. پژوهش \cite{neel2021descent} الگوریتم‌های مبتنی بر کاهش شیب با تضمین ارائه می‌دهد و \cite{sekhari2021remember} موازنه‌ی میان تعمیم و حذف را بررسی می‌کند.

باید صریح گفت که این تضمین‌ها به تنظیم این پایان‌نامه منتقل نمی‌شوند. آموزش دنباله‌ای یک شبکه‌ی عمیق غیرمحدب، با نقاب‌های وابسته به تاریخ و یک حافظه‌ی بازپخش که خودش تابعی از دنباله‌ی درخواست‌ها است، هیچ‌یک از پیش‌فرض‌های آن قضایا را برآورده نمی‌کند. به همین دلیل، همه‌ی ادعاهای این پایان‌نامه تجربی برچسب می‌خورند. بحث سازوکارهای «یادگیری برای فراموش‌کردن» در \cite{qu2023learntounlearn} آمده است.

برای آنکه روشن باشد چه چیزی منتقل نمی‌شود، صورت آرام‌شده‌ی این تضمین را می‌آوریم. الگوریتم حذف $U$ را نسبت به الگوریتم آموزش $A$ و مجموعه‌داده‌ی $D$ یک حذف $(\epsilon,\delta)$ می‌نامیم هرگاه برای هر نمونه‌ی $z \in D$ و هر مجموعه‌ی اندازه‌پذیر $S$ از مدل‌ها، رابطه‌ی
\begin{equation}
\label{eq:certified-unlearning-guarantee}
\Pr\big[U(A(D), z) \in S\big] \le e^{\epsilon}\,\Pr\big[A(D \setminus z) \in S\big] + \delta
\end{equation}
و رابطه‌ی متقارن آن برقرار باشد. این صورت‌بندی عمدا به تعریف حریم خصوصی تفاضلی \cite{dwork2014algorithmic} شبیه است و مانند آن، گزاره‌ای درباره‌ی \emph{توزیع} خروجی الگوریتم است و نه درباره‌ی یک اجرای منفرد.

سه پیش‌فرض این قضایا در تنظیم ما نقض می‌شوند. نخست، اثبات‌ها به محدب‌بودن تابع هزینه و یکتایی کمینه تکیه دارند تا بتوانند فاصله‌ی پیراسنجه‌ها را کران‌دار کنند؛ یک شبکه‌ی عمیق هیچ‌کدام را ندارد. دوم، فرض می‌شود مسیر آموزش تنها از راه داده به آن وابسته است، در حالی که اینجا نقاب‌های زیرشبکه خودشان تابعی از ترتیب وظایف‌اند و تاریخ اجرا را در ساختار مدل رمز می‌کنند. سوم، حافظه‌ی بازپخش محدود، خود تابعی از دنباله‌ی درخواست‌هاست؛ پس حتی اگر به‌روزرسانی حذف کران‌دار می‌بود، وضعیت اولیه‌ای که روی آن اعمال می‌شود به داده‌ی حذف‌شونده وابسته است. به همین سه دلیل، در این پایان‌نامه هیچ عدد $\epsilon$ گزارش نمی‌شود و هر ادعای حذف، تجربی و مشروط بیان می‌گردد.

\subsection{یادگیری‌زدایی در بستر یادگیری پیوسته}
\label{subsec:lit-cl-unlearning}

تلاقی این دو حوزه، موضوع پژوهش حاضر است و چند مرجع مهم دارد. نخست، \lr{CLPU} \cite{liu2022clpu} یک مرجع مبتنی بر جداسازی کامل است. در پیاده‌سازی ارزیابی‌شده، هر وظیفه‌ی فعال شبکه‌ی کامل خود را دارد و با دریافت درخواست حذف، شبکه‌ی مربوط به همان وظیفه کنار گذاشته می‌شود. در نتیجه، تداخل میان وظیفه‌ها از نظر ساختاری رخ نمی‌دهد و بیشترین افت همواره صفر است. هزینه‌ی این جداسازی، افزایش حافظه‌ی موردنیاز با تعداد وظیفه‌ها است. این پایان‌نامه \lr{CLPU} را به‌عنوان مرجع جداسازی کامل در نظر می‌گیرد و در بخش \ref{sec:res-storage} حافظه‌ی موردنیاز آن را مستقیما محاسبه می‌کند.

\lr{PALL} \cite{ozdenizci2025pall} نقطه‌ی شروع فنی این کار است. \lr{PALL} برای هر وظیفه یک زیرشبکه‌ی تنک می‌آموزد، نقاب آن را ذخیره می‌کند و در زمان درخواست، بر اساس اینکه یک درایه در نقاب‌های تاریخی و پشتیبان‌های ذخیره‌شده چه وضعیتی دارد، آن را بازنشانی یا بازسازی می‌کند و سپس وظایف نگه‌داری‌شده را ترمیم می‌کند. آنچه \lr{PALL} فراهم می‌کند، ساختاری است که هم‌پوشانی را قابل‌اندازه‌گیری می‌سازد؛ آنچه فراهم نمی‌کند، تفکیک درایه‌های مشترک بر حسب حساسیت نگه‌داری است.

نزدیک‌ترین کار به صورت‌مسئله‌ی ما {یادگیری با فراموشی انتخابی\en{Learning with Selective Forgetting (LSF)}} \cite{shibata2021lsf} است که یک کلاس را با کدهای یادیار پاک می‌کند در حالی که یادگیری ادامه دارد؛ همین جفت‌شدگی یادگیری و حذف در راهبرد پادگرا برای فراموشی انتخابی مادام‌العمر \cite{shan2024contrastive} نیز تکرار می‌شود. تفاوت تنظیم ما با هر دو در آن است که یک مدل مشترک با حافظه‌ی بازپخش کران‌دار نگه می‌داریم و وزن‌های هم‌پوشان را مستقیما ویرایش می‌کنیم، نه آنکه فراموشی را از راه هدف یادگیری اعمال کنیم.

دو پژوهش جدیدتر دامنه‌ی ارزیابی را گسترش می‌دهند. \lr{UnCLe} \cite{adhikari2026uncle} پیراسنجه‌های مخصوص هر وظیفه را تولید می‌کند و {بازگشت وظیفه پس از حذف\en{Post-Deletion Task Relapse}} را بدون داده‌های نگه‌داری‌شده بررسی می‌کند. \lr{CATA} \cite{lin2026cata} حذف دنباله‌ای و ماندگاری آن را برای مدل‌های بینایی‌ـ‌زبان، از طریق تجمیع بردار وظیفه و پرهیز از تعارض، مطالعه می‌کند. پژوهش \cite{chatterjee2024unified} نیز نشان می‌دهد تقطیر کنترل‌شده می‌تواند یادگیری و حذف را با حافظه‌ی ثابت هم‌زمان پشتیبانی کند. درس مشترک این پژوهش‌ها آن است که دقت بلافاصله پس از درخواست، شاهد کافی نیست و ماندگاری سرکوب باید جداگانه ارزیابی شود. بخش \ref{subsec:res-persistence} همین ارزیابی را انجام می‌دهد.

\section{یادگیری کم‌پیراسنجه}
\label{sec:lit-peft}

{یادگیری کم‌پیراسنجه\en{Parameter-Efficient Fine-Tuning (PEFT)}} آموزش را به پیمانه‌های فشرده‌ی درج‌شده در یک شبکه‌ی منجمد محدود می‌کند. {لایه‌های تطبیق گلوگاهی\en{Bottleneck Adapters}} \cite{houlsby2019parameter} یک نگاشت کاهش ابعاد، یک غیرخطیت \lr{ReLU} \cite{nair2010relu} و یک نگاشت افزایش ابعاد را در قالب یک بلوک باقی‌مانده ترکیب می‌کنند:
\begin{equation}
	\label{eq:adapter}
	h \;\leftarrow\; h + W_{up}\,\sigma\!\big(W_{down}\,h\big),
	\qquad W_{down}\in\mathbb{R}^{r\times d},\ W_{up}\in\mathbb{R}^{d\times r},\ r\ll d.
\end{equation}
{انطباق کم‌رتبه\en{Low-Rank Adaptation}} یا \lr{LoRA} \cite{hu2021lora} به‌جای افزودن یک پیمانه، به‌روزرسانی وزن را با حاصل‌ضرب دو ماتریس کم‌رتبه‌ی $\Delta W=BA$ بیان می‌کند. لایه‌های تطبیق باقی‌مانده‌ی چندحوزه‌ای \cite{rebuffi2017residualadapters}، زیرساخت‌های اشتراک لایه‌ی تطبیق \cite{pfeiffer2020adapterhub}، نسخه‌های فشرده‌تر \cite{mahabadi2021compacter,ruckle2021adapterdrop} و روش‌های مبتنی بر پیشوند و راهنمای ورودی \cite{li2021prefixtuning,lester2021prompttuning} این خانواده را تکمیل می‌کنند. پیش‌نویس اخیر \lr{BID-LoRA} \cite{rachapudi2026bidlora} ترکیب مشابهی از یادگیری پیوسته و یادگیری‌زدایی را بررسی می‌کند.

بررسی این خانواده در پایان‌نامه دو دلیل دارد. نخست، اگر بدنه‌ی اصلی ثابت باشد، هنگام حذف فقط بخش کوچکی از پیراسنجه‌ها تغییر می‌کنند. دوم، در یک لایه‌ی تطبیق مشترک و باریک، هم‌پوشانی میان وظیفه‌ها متمرکز است؛ بنابراین انتظار می‌رود سازوکار آگاه از هم‌پوشانی در این بخش اثر روشنی داشته باشد. فصل \ref{ch:results} نشان می‌دهد این انتظار برآورده نمی‌شود و نتیجه‌ی منفی نیز در گزارش حفظ شده است.

\section{ارزیابی فراموشی و حریم خصوصی}
\label{sec:lit-eval}

\subsection{معیارهای رفتاری}
\label{subsec:lit-metrics}

سنجش کیفیت فراموشی معمولا با معیارهای رفتاری آغاز می‌شود: دقت روی وظیفه‌ی هدف پس از درخواست، دقت روی وظایف نگه‌داری‌شده و اختلاف این دو نسبت به یک مرجع. نکته‌ی ظریفی که در این معیارها اغلب نادیده گرفته می‌شود، جهت مطلوب دقت وظیفه‌ی حذف‌شده است. کاهش این دقت تا صفر مطلوب نیست، چون مدلی که روی وظیفه‌ی حذف‌شده \emph{به‌طور نظام‌مند} اشتباه می‌کند، همچنان اطلاعاتی از آن وظیفه حمل می‌کند و یک حمله‌کننده می‌تواند از همین سوگیری برای تشخیص عضویت استفاده کند. هدف درست، رساندن دقت به سطح تصادفی است؛ یعنی مدل باید نسبت به وظیفه‌ی حذف‌شده \emph{بی‌اطلاع} به نظر برسد و نه بدرفتار. به همین دلیل در این پایان‌نامه هرجا دقت وظیفه‌ی هدف گزارش می‌شود، معیار داوری فاصله‌ی آن تا سطح تصادفی است و نه مقدار خام آن. تمام جدول‌های فصل \ref{ch:results} بر همین اساس تفسیر می‌شوند.

\subsection{حمله‌ی استنتاج عضویت}
\label{subsec:lit-mia}

{حمله‌ی استنتاج عضویت\en{Membership Inference Attack (MIA)}} \cite{shokri2017membership,yeom2018privacy} تلاش می‌کند تشخیص دهد آیا یک نمونه‌ی مشخص در آموزش مدل حضور داشته است یا نه؛ ازاین‌رو، به‌عنوان یک آزمون تجربی حریم خصوصی به‌کار می‌رود. نسخه‌ی ساده‌ی این حمله از آستانه‌ی زیان استفاده می‌کند. دو نکته درباره‌ی این معیار مهم است. نخست، پیش‌پردازش نمونه‌های عضو و غیرعضو باید یکسان باشد؛ در غیر این صورت، تفاوت مشاهده‌شده ممکن است فقط از پیش‌پردازش متفاوت ناشی شود. دوم، حمله‌ی مبتنی بر آستانه‌ی زیان ممکن است توان تشخیص کمی داشته باشد. حمله‌های قوی‌تر مبتنی بر مدل‌های سایه \cite{carlini2022membership} و ارزیابی‌های جامع‌تر \cite{nasr2019comprehensive} می‌توانند تفاوتی را آشکار کنند که آزمون ساده تشخیص نمی‌دهد. بنابراین، نتیجه‌ی نزدیک به $0.5$ فقط نشان می‌دهد حمله‌ی به‌کاررفته تفاوتی پیدا نکرده است و نشانه‌ی حذف موفق نیست.

خانواده‌ی حمله‌های دیگر نیز زمینه‌ی تهدید را کامل می‌کنند: {وارون‌سازی مدل\en{Model Inversion}} \cite{fredrikson2015model} و استخراج داده‌ی آموزشی از مدل‌های زبانی \cite{carlini2021extracting} نشان می‌دهند اثر داده در وزن‌ها تا چه اندازه قابل بازیابی است.

\subsection{مرجع بازآموزی و ارزیابی رسمی}
\label{subsec:lit-retrain}

یک ارزیابی دقیق‌تر، مقایسه‌ی مدل فراموش‌کرده با یک نمونه‌ی تازه از همان روش است که از ابتدا بدون وظیفه‌ی هدف آموزش دیده باشد. سنجه‌های مقایسه شامل توافق پیش‌بینی، فاصله‌ی {جنسن-شانون\en{Jensen--Shannon}} میان توزیع‌های خروجی، فاصله‌ی $\ell_2$ امتیازهای خام و شباهت کسینوسی ویژگی‌ها است. در سوی صوری، \cite{steinke2023privacy} روشی برای ارزیابی حریم خصوصی با یک اجرای آموزش ارائه می‌دهد و {حریم خصوصی تفاضلی\en{Differential Privacy}} \cite{dwork2014algorithmic} چارچوب نظری استاندارد این حوزه است. بخش \ref{subsec:res-retrain} مرجع بازآموزی هم‌روش را برای \lr{EPALL} گزارش می‌کند: نتیجه شاهدی برای هم‌ارزی با بازآموزی دقیق فراهم نمی‌کند، هرچند بدون یک خط پایه‌ی «بازآموزی در برابر بازآموزی» نمی‌توان از آن، حذف دقیق توزیعی را به‌طور صوری رد کرد.

شکل \ref{fig:lit-taxonomy} مسیر مفهومی ادبیات را از حفظ دانش گذشته تا حذف انتخابی آگاه از هم‌پوشانی خلاصه می‌کند. این شکل یک رده‌بندی مفهومی است و نه ترتیب تاریخی یا مقایسه‌ی تجربی.

\begin{figure}[htbp]
	\centering
	\includegraphics[width=0.98\columnwidth]{literature_taxonomy.png}
	\caption[رده‌بندی مفهومی روش‌های مرتبط]{رده‌بندی مفهومی روش‌های مرتبط و جایگاه این پایان‌نامه. مسیر از روش‌های حفظ دانش در یادگیری پیوسته و جداسازی پیراسنجه‌ها، به یادگیری‌زدایی و سپس حذف وظیفه‌ای آگاه از هم‌پوشانی می‌رسد. این رده‌بندی برای سازمان‌دهی مفاهیم است و ادعای تجربی یا ترتیب تاریخی ندارد.}
	\label{fig:lit-taxonomy}
\end{figure}

\section{مقایسه‌ی مفهومی و جایگاه این پژوهش}
\label{sec:lit-gap}

جدول \ref{tab:lit-comparison} روش‌های مرتبط را از چهار جنبه مقایسه می‌کند: مشترک‌بودن مدل میان وظیفه‌ها، پشتیبانی از حذف در سطح وظیفه، مدیریت مستقیم هم‌پوشانی پیراسنجه‌ها و تغییر حافظه‌ی موردنیاز با افزایش تعداد وظیفه‌ها.

\begin{table}[htbp]
    \centering
    \footnotesize
    \caption[مقایسه‌ی مفهومی روش‌های مرتبط]{مقایسه‌ی مفهومی نزدیک‌ترین روش‌های مرتبط بر چهار محور تصمیم‌گیری این پایان‌نامه. دو سطر پایانی روش‌های پیشنهادی این کار هستند.}
    \label{tab:lit-comparison}
    \setlength{\tabcolsep}{3pt}
    \renewcommand{\arraystretch}{0.94}
    \begin{adjustbox}{max width=\textwidth,max totalheight=0.82\textheight,keepaspectratio}
    \begin{tabular}{p{3.3cm} p{3.0cm} p{1.6cm} p{4.4cm} p{3.1cm}}
    \toprule
    \textbf{روش} & \textbf{مدل مشترک} & \textbf{حذف سطح وظیفه} & \textbf{مدیریت مستقیم هم‌پوشانی} & \textbf{رشد حافظه‌ی موردنیاز} \\ \midrule
    \lr{EWC}~\cite{kirkpatrick2017overcoming} & بله & خیر & خیر (جریمه‌ی سراسری) & ثابت \\
    \lr{LwF}~\cite{li2017learning} & بله & خیر & خیر (تقطیر خروجی) & ثابت \\
    \lr{ER / DER++}~\cite{buzzega2020dark} & بله & خیر & خیر & فقط حافظه \\
    \mbox{\lr{GEM / A-GEM}}~\cite{lopezpaz2017gem} & بله & خیر & ضمنی، تصویر سراسری & فقط حافظه \\
    \lr{PackNet}~\cite{mallya2018packnet} & بله & خیر & نقاب دارد، وزن منجمد & نقاب \\
    \lr{WSN}~\cite{kang2022wsn} & بله & خیر & نقاب دارد، برای حذف به‌کار نمی‌رود & نقاب \\
    \lr{SSD}~\cite{foster2024ssd} & بله & بله & فقط میراکردن اهمیتی & ثابت \\
    \lr{SalUn}~\cite{fan2024salun} & بله & بله & انتخاب وزن برای \emph{تغییر} & ثابت \\
    \lr{CLPU}~\cite{liu2022clpu} & خیر & بله & به‌طور ساختاری لازم نیست & یک شبکه به‌ازای وظیفه \\
    \lr{PALL}~\cite{ozdenizci2025pall} & بله & بله & نقاب ذخیره، مشترک‌ها بی‌رتبه & نقاب و پشتیبان \\
    \addlinespace[3pt]
    \textbf{\lr{EPALL}} & بله & بله & \textbf{محافظت رتبه‌بندی‌شده‌ی مشترک} & نقاب و پشتیبان \\
    \textbf{\mbox{\lr{PALL-Adapter}}} & بله & بله & نقاب نرم (ارزیابی‌شده، بی‌اثر) & لایه‌ی تطبیق به‌ازای وظیفه \\
    \bottomrule
    \end{tabular}
    \end{adjustbox}
\end{table}


تفسیر این جدول، خلأ پژوهشی را روشن می‌کند. روش‌های تنظیم و بازپخش، سازوکار حذف در سطح وظیفه ندارند. روش‌های یادگیری‌زدایی عمومی مانند \lr{SSD} و \lr{SalUn} سازوکار حذف دارند، اما ساختار وظیفه‌ای مدل را در نظر نمی‌گیرند و، چنان‌که فصل \ref{ch:results} نشان می‌دهد، در محیط پیوسته دقت وظیفه‌ی هدف را به سطح حدس تصادفی نمی‌رسانند. \lr{CLPU} با جداسازی کامل، هم‌پوشانی را از بین می‌برد؛ بهای این جداسازی، رشد خطی حافظه‌ی موردنیاز است. روش‌هایی مانند \lr{WSN} و \lr{PackNet} نقاب‌های مربوط به هر وظیفه را می‌سازند، اما از آن‌ها برای حذف استفاده نمی‌کنند. \lr{PALL} این دو ایده را به هم پیوند می‌دهد، ولی وزن‌های مشترک را بر پایه‌ی اهمیتشان رتبه‌بندی نمی‌کند.

این پایان‌نامه خلأ باقی‌مانده را در یک سامانه‌ی دارای مدل مشترک بررسی می‌کند. با استفاده از نقاب هر وظیفه، هم‌پوشانی ساختاری محاسبه می‌شود؛ سپس وزن‌های مشترک بر پایه‌ی حساسیت وظیفه‌های نگه‌داری‌شده رتبه‌بندی می‌شوند و فقط زیرمجموعه‌ی مهم آن‌ها هنگام ترمیم محافظت می‌شود. فصل بعد این ایده را به یک الگوریتم مشخص تبدیل می‌کند.
